\section{Marcas}
\label{sec:marcas}

Las marcas son elementos gráficos que se colocan en las intersecciones del tablero para resaltar posiciones específicas, indicar puntos de interés o numerar secuencias, sin necesidad de añadir piedras. \pgfgoname\ ofrece comandos y opciones versátiles para personalizar estas marcas, desde formas predefinidas hasta etiquetas personalizadas.

\begin{command}{\mark\opt{\oarg{options}}\marg{intersection}}
	
	Dibuja una marca individual en la intersección especificada (ej. «C3»). Las opciones permiten definir su forma, color y otros atributos.
\end{command}

\begin{command}{\marks\opt{\oarg{options}}\marg{intersection \opt{, intersection, ...}}}
	
	Similar a |\mark|, pero aplica la misma marca a múltiples intersecciones mediante la sintaxis de coordenadas.
\end{command}

\begin{key}{/pgfgo/mark = \meta{options} (default cross)}
	
	Define el tipo de marca a colocar, con una cruz («cross») como valor predeterminado si no se especifica otra opción.
	
	\begin{codeexample}[]
		\pgfgoset{goban = {5, partial = {from = B2 to D4}}}
		
		\begin{goban}
			\stones[b]{C2,D3}
			\stone[w]{B4}
		\end{goban}
		
		\begin{goban}
			\stones[b]{C2,D3}
			\stone[w]{B4}
			\mark{C3}
		\end{goban}
	\end{codeexample}	
\end{key}

\begin{key}{/pgfgo/mark/circle}
	
	Coloca un círculo hueco como marca.
	
	\begin{codeexample}[]
		\pgfgoset{goban = {5, partial = {from = B2 to D4}}}
		
		\begin{goban}
			\stones[b]{C2,D3}
			\stone[w]{B4}
		\end{goban}
		
		\begin{goban}
			\stones[b]{C2,D3}
			\stone[w]{B4}
			\mark[circle]{C3}
		\end{goban}
	\end{codeexample}	
\end{key}

\begin{key}{/pgfgo/mark/triangle}
	
	Dibuja un triángulo hueco.
	
	\begin{codeexample}[]
		\pgfgoset{goban = {5, partial = {from = B2 to D4}}}
		
		\begin{goban}
			\stones[b]{C2,D3}
			\stone[w]{B4}
		\end{goban}
		
		\begin{goban}
			\stones[b]{C2,D3}
			\stone[w]{B4}
			\mark[triangle]{C3}
		\end{goban}
	\end{codeexample}	
\end{key}

\begin{key}{/pgfgo/mark/square}
	
	Coloca un cuadrado hueco.
	
	\begin{codeexample}[]
		\pgfgoset{goban = {5, partial = {from = B2 to D4}}}
		
		\begin{goban}
			\stones[b]{C2,D3}
			\stone[w]{B4}
		\end{goban}
		
		\begin{goban}
			\stones[b]{C2,D3}
			\stone[w]{B4}
			\mark[square]{C3}
		\end{goban}
	\end{codeexample}
\end{key}

\begin{key}{/pgfgo/mark/cross}
	
	Dibuja una cruz.
	
	\begin{codeexample}[]
		\pgfgoset{goban = {5, partial = {from = B2 to D4}}}
		
		\begin{goban}
			\stones[b]{C2,D3}
			\stone[w]{B4}
		\end{goban}
		
		\begin{goban}
			\stones[b]{C2,D3}
			\stone[w]{B4}
			\mark[cross]{C3}
		\end{goban}
	\end{codeexample}
\end{key}

\begin{key}{/pgfgo/mark/filled square}
	
	Coloca un cuadrado relleno.
	
	\begin{codeexample}[]
		\pgfgoset{goban = {5, partial = {from = B2 to D4}}}
		
		\begin{goban}
			\stones[b]{C2,D3}
			\stone[w]{B4}
		\end{goban}
		
		\begin{goban}
			\stones[b]{C2,D3}
			\stone[w]{B4}
			\mark[filled square]{C3}
		\end{goban}
	\end{codeexample}
\end{key}

\begin{key}{/pgfgo/mark/filled circle}
	
	Dibuja un círculo relleno.
	
	\begin{codeexample}[]
		\pgfgoset{goban = {5, partial = {from = B2 to D4}}}
		
		\begin{goban}
			\stones[b]{C2,D3}
			\stone[w]{B4}
		\end{goban}
		
		\begin{goban}
			\stones[b]{C2,D3}
			\stone[w]{B4}
			\mark[filled circle]{C3}
		\end{goban}
	\end{codeexample}
\end{key}

\begin{key}{/pgfgo/mark/t-circle}
	
	Coloca un círculo pequeño con un borde más grueso, diseñado específicamente para marcar territorios. También puede escribirse |tcircle|.
	
	\begin{codeexample}[]
		\pgfgoset{goban = {5, partial = {from = B2 to D4}}}
		
		\begin{goban}
			\stones[b]{C2,D3}
			\stone[w]{B4}
		\end{goban}
		
		\begin{goban}
			\stones[b]{C2,D3}
			\stone[w]{B4}
			\mark[t-circle]{C3}
		\end{goban}
	\end{codeexample}
\end{key}

\begin{key}{/pgfgo/mark/t-square}
	
	Dibuja un cuadrado pequeño con borde grueso, similar a «t-circle», usado para señalar territorios. También puede escribirse |tsquare|.
	
	\begin{codeexample}[]
		\pgfgoset{goban = {5, partial = {from = B2 to D4}}}
		
		\begin{goban}
			\stones[b]{C2,D3}
			\stone[w]{B4}
		\end{goban}
		
		\begin{goban}
			\stones[b]{C2,D3}
			\stone[w]{B4}
			\mark[t-square]{C3}
		\end{goban}
	\end{codeexample}
\end{key}

\begin{key}{/pgfgo/mark/hatched lines}
	
	Coloca un patrón de líneas cruzadas (tramado).
	
	\begin{codeexample}[]
		\pgfgoset{goban = {5, partial = {from = B2 to D4}}}
		
		\begin{goban}
			\stones[b]{C2,D3}
			\stone[w]{B4}
		\end{goban}
		
		\begin{goban}
			\stones[b]{C2,D3}
			\stone[w]{B4}
			\mark[hatched lines]{C3}
		\end{goban}
	\end{codeexample}
\end{key}

\begin{key}{/pgfgo/mark/sequence=\meta{options} (default from = 1)}
	Permite numerar una serie de marcas en orden.
\end{key}

\begin{key}{/pgfgo/mark/sequence/use}
	Muestra la marca |sequence| sobre el tablero. Es lo mismo que |sequence = true|.
\end{key}

\begin{key}{/pgfgo/mark/sequence/false}
	Desactiva la secuencia (por ejemplo, una activada con |\pgfgoset|).
\end{key}

\begin{key}{/pgfgo/mark/sequence/format = \meta{format}}
	Activa la secuencia y elige el formato de los números. Es lo mismo que |sequence, label = {format = ...}| (ver |label/format|).
\end{key}

\begin{key}{/pgfgo/mark/sequence/from=\meta{number}}
	
	Establece el número inicial de la secuencia de marcas numeradas, permitiendo continuar desde un punto específico en lugar de 1.
	
	\begin{codeexample}[]
		\pgfgoset{goban = {5, partial = {from = B2 to D4}}}
		
		\begin{goban}
			\stones[b]{C2,D3}
			\stone[w]{B4}
		\end{goban}
		
		\begin{goban}
			\stones[b]{C2,D3}
			\stone[w]{B4}
			\marks[sequence = {use, from = 3}]{C3,B3,C4}
		\end{goban}
	\end{codeexample}
\end{key}




\begin{key}{/pgfgo/mark/dim}
	Atenúa el punto: lo cubre con el color del fondo del tablero, translúcido, sobre la piedra y la cuadrícula. Es lo que dibuja |\sgf| para la propiedad |DD|.

	\begin{codeexample}[]
		\begin{goban}[5]
			\stones{B2, C3}\stone[w]{D4}
			\marks[dim]{A1:5, B2}
		\end{goban}
	\end{codeexample}
\end{key}

\begin{key}{/pgfgo/mark/custom=\meta{token list}}
	
	Permite definir una marca personalizada usando texto o símbolos (como «!» o «?») e incluso im\'agenes.
	
	\begin{codeexample}[]
		\pgfgoset{goban = {5, partial = {from = B2 to D4}}}
		
		\begin{goban}
			\stones[b]{C2,D3}
			\stone[w]{B4}
		\end{goban}
		
		\begin{goban}
			\stones[b]{C2,D3}
			\stone[w]{B4}
			\mark[custom = {!}]{C3}
		\end{goban}
		
		\begin{goban}
			\stones[b]{C2,D3}
			\stone[w]{B4}
			\mark[custom = {\tikz{\node[rounded corners, draw, inner sep = 1em, red]{\dots};}}]{C3}
		\end{goban}
	\end{codeexample}	
\end{key}


\begin{key}{/pgfgo/mark/line = \meta{options}}
	Controla el contorno de la marca.
\end{key}


\begin{key}{/pgfgo/mark/line/color = \meta{color}}
	Define el color del contorno de la marca.
	
	\begin{codeexample}[]
		\pgfgoset{goban = {5, partial = {from = B2 to D4}}}
		
		\begin{goban}
			\stones[b]{C2,D3}
			\stone[w]{B4}
		\end{goban}
		
		\begin{goban}
			\stones[b]{C2,D3}
			\stone[w]{B4}
			\mark[triangle, line = {color = red}]{C3}
		\end{goban}
	\end{codeexample}
\end{key}


\begin{key}{/pgfgo/mark/line/width = \meta{dimension}}
	
	Ajusta el grosor del contorno de las marcas.
	
	\begin{codeexample}[]
		\pgfgoset{goban = {5, partial = {from = B2 to D4}}}
		
		\begin{goban}
			\stones[b]{C2,D3}
			\stone[w]{B4}
			\mark[triangle]{C3}
		\end{goban}
		
		\begin{goban}
			\stones[b]{C2,D3}
			\stone[w]{B4}
			\mark[triangle, line = {width = 0.1pt}]{C3}
		\end{goban}
	\end{codeexample}
\end{key}



\begin{key}{/pgfgo/mark/fill = \meta{options}}
	Controla el relleno de las marcas rellenas (|filled square|, |filled circle| y |hatched lines|).
\end{key}

\begin{key}{/pgfgo/mark/fill/true}
	Las marcas rellenas usan |fill/color|. Sin |fill = true|, se rellenan de gris.
\end{key}

\begin{key}{/pgfgo/mark/fill/false}
	Las marcas rellenas vuelven a rellenarse de gris. Es el valor inicial.
\end{key}

\begin{key}{/pgfgo/mark/fill/color = \meta{color}}
	
	Establece el color de relleno de una marca. Solo se usa con |fill = true|.
	
	\begin{codeexample}[]
		\pgfgoset{goban = {5, partial = {from = B2 to D4}}}
		
		\begin{goban}
			\stones[b]{C2,D3}
			\stone[w]{B4}
		\end{goban}
		
		\begin{goban}
			\stones[b]{C2,D3}
			\stone[w]{B4}
			\mark[filled circle, fill = {true, color = red}]{C3}
		\end{goban}
	\end{codeexample}
\end{key}


\begin{key}{/pgfgo/mark/label = \meta{option}}
	Controla la apariencia de las etiquetas de las marcas que las tengan.
\end{key}


\begin{key}{/pgfgo/mark/label/color = \meta{color}}
	
	Controla el color del texto o símbolo en marcas personalizadas (como con |custom|).
	
	\begin{codeexample}[]
		\pgfgoset{goban = {5, partial = {from = B2 to D4}}}
		
		\begin{goban}
			\stones[b]{C2,D3}
			\stone[w]{B4}
		\end{goban}
		
		\begin{goban}
			\stones[b]{C2,D3}
			\stone[w]{B4}
			\mark[custom = {!}, label = {color = red}]{C3}
		\end{goban}
	\end{codeexample}
\end{key}

\begin{key}{/pgfgo/mark/label/font = \meta{font commands}}
	Fuente de las etiquetas de las marcas (|sequence| y |custom|).

	\begin{codeexample}[]
		\pgfgoset{goban = {5, partial = {from = B2 to D4}}}

		\begin{goban}
			\stones[b]{C2,D3}
			\stone[w]{B4}
			\marks[sequence, label = {font = \itshape}]{C3,B3,C4}
		\end{goban}
	\end{codeexample}
\end{key}

\begin{key}{/pgfgo/mark/label/format=\meta{options} (initially arabic)}
	
	Controla el estilo de numeración de las etiquetas de enumeraci\'on (como |sequence|), con opciones como números arábigos, letras o romanos.
\end{key}

\begin{key}{/pgfgo/mark/label/format/arabic}
	
	Usa números arábigos para secuencias.
	
	\begin{codeexample}[]
		\pgfgoset{goban = {5, partial = {from = B2 to D4}}}
		
		\begin{goban}
			\stones[b]{C2,D3}
			\stone[w]{B4}
		\end{goban}
		
		\begin{goban}
			\stones[b]{C2,D3}
			\stone[w]{B4}
			\marks[label = {format = arabic}, sequence]{C3,B3,C4}
		\end{goban}
	\end{codeexample}
\end{key}

\begin{key}{/pgfgo/mark/label/format/alph}
	
	Usa letras minúsculas para la secuencia.
	
	\begin{codeexample}[]
		\pgfgoset{goban = {5, partial = {from = B2 to D4}}}
		
		\begin{goban}
			\stones[b]{C2,D3}
			\stone[w]{B4}
		\end{goban}
		
		\begin{goban}
			\stones[b]{C2,D3}
			\stone[w]{B4}
			\marks[label = {format = alph}, sequence]{C3,B3,C4}
		\end{goban}
	\end{codeexample}
\end{key}

\begin{key}{/pgfgo/mark/label/format/Alph}
	
	Usa letras mayúsculas.
	
	\begin{codeexample}[]
		\pgfgoset{goban = {5, partial = {from = B2 to D4}}}
		
		\begin{goban}
			\stones[b]{C2,D3}
			\stone[w]{B4}
		\end{goban}
		
		\begin{goban}
			\stones[b]{C2,D3}
			\stone[w]{B4}
			\marks[label = {format = Alph}, sequence]{C3,B3,C4}
		\end{goban}
	\end{codeexample}
\end{key}

\begin{key}{/pgfgo/mark/label/format/roman}
	
	Usa números romanos minúsculos.
	
	\begin{codeexample}[]
		\pgfgoset{goban = {5, partial = {from = B2 to D4}}}
		
		\begin{goban}
			\stones[b]{C2,D3}
			\stone[w]{B4}
		\end{goban}
		
		\begin{goban}
			\stones[b]{C2,D3}
			\stone[w]{B4}
			\marks[label = {format = roman}, sequence]{C3,B3,C4}
		\end{goban}
	\end{codeexample}
\end{key}

\begin{key}{/pgfgo/mark/label/format/Roman}
	
	Usa números romanos mayúsculos.
	
	\begin{codeexample}[]
		\pgfgoset{goban = {5, partial = {from = B2 to D4}}}
		
		\begin{goban}
			\stones[b]{C2,D3}
			\stone[w]{B4}
		\end{goban}
		
		\begin{goban}
			\stones[b]{C2,D3}
			\stone[w]{B4}
			\marks[label = {format = Roman}, sequence]{C3,B3,C4}
		\end{goban}
	\end{codeexample}
\end{key}



\begin{key}{/pgfgo/mark/black}
	
	Aplica las propiedades del primer jugador (negro) a la marca.
	
	\begin{codeexample}[]
		\pgfgoset{goban = {5, partial = {from = B2 to D4}}}
		
		\begin{goban}
			\stones[b]{C2,D3}
			\stone[w]{B4}
			\mark[triangle]{C3}
		\end{goban}
		
		\begin{goban}
			\stones[b]{C2,D3}
			\stone[w]{B4}
			\mark[black, t-square]{C3}
		\end{goban}
	\end{codeexample}
\end{key}

\begin{key}{/pgfgo/mark/white}
	
	Aplica las propiedades del segundo jugador (blanco) a la marca.
\end{key}

\begin{key}{/pgfgo/mark/neutral}
	
	Mantiene la marca sin asociarla a ningún jugador.
\end{key}

\begin{key}{/pgfgo/mark/b}
	Abreviatura de |black|.
\end{key}

\begin{key}{/pgfgo/mark/w}
	Abreviatura de |white|.
\end{key}

\begin{key}{/pgfgo/mark/n}
	Abreviatura de |neutral|.
\end{key}

\begin{key}{/pgfgo/mark/forget}
	Esta marca no se guarda en la memoria del tablero, aunque el |goban| use |remember|.
\end{key}

\begin{key}{/pgfgo/mark/remember}
	Esta marca se guarda en la memoria del tablero, aunque el |goban| use |forget|.
\end{key}


\subsection{Flechas y líneas}
\label{subsec:flechas}

Las flechas y las líneas unen dos intersecciones. Se dibujan encima de las piedras, con un halo del color del fondo del tablero que las separa de las piedras y de la cuadrícula. Como |\mark|, |\arrow| y |\line| solo existen dentro de |goban|: afuera, |\line| sigue siendo el comando de \LaTeX. |\sgf| dibuja las flechas (|AR|) y las líneas (|LN|) del archivo.

\begin{command}{\arrow\opt{\oarg{options}}\marg{from}\marg{to}}
	Dibuja una flecha de una intersección a otra.

	\begin{codeexample}[]
		\begin{goban}[9]
			\stones{C3, D3}\stone[w]{E5}
			\arrow{B3}{F3}
			\arrow[color = red, halo = false]{C7}{G9}
		\end{goban}
	\end{codeexample}
\end{command}

\begin{command}{\line\opt{\oarg{options}}\marg{from}\marg{to}}
	Dibuja una línea de una intersección a otra.

	\begin{codeexample}[]
		\begin{goban}[9]
			\stones{C3, D4}\stone[w]{E5}
			\line{C3}{G7}
			\line[color = blue, width = 2pt]{A9}{J9}
		\end{goban}
	\end{codeexample}
\end{command}

\begin{key}{/pgfgo/arrow/color=\meta{color} (initially black)}
	Color de la flecha.
\end{key}

\begin{key}{/pgfgo/arrow/width=\meta{dimension} (initially 1.2pt)}
	Grosor de la flecha. La punta crece con él.
\end{key}

\begin{key}{/pgfgo/arrow/opacity=\meta{number} (initially 1)}
	Opacidad de la flecha.
\end{key}

\begin{key}{/pgfgo/arrow/tip=\meta{arrow tip} (initially Stealth)}
	La punta de la flecha, una de las de la biblioteca |arrows.meta| de \tikzname, como |Latex| o |Triangle|.
\end{key}

\begin{key}{/pgfgo/arrow/halo=\meta{halo options} (default true)}
	El halo alrededor de la flecha.
\end{key}

\begin{stylekey}{/pgfgo/arrow/halo/true}
	Dibuja el halo. Es el valor inicial.
\end{stylekey}

\begin{stylekey}{/pgfgo/arrow/halo/false}
	No dibuja el halo.
\end{stylekey}

\begin{key}{/pgfgo/arrow/halo/color=\meta{color}}
	Color del halo. Sin esta \emph{key}, es el del fondo del tablero.
\end{key}

\begin{key}{/pgfgo/arrow/halo/width=\meta{dimension} (initially 0.8pt)}
	Ancho del halo a cada lado de la flecha.
\end{key}

\begin{key}{/pgfgo/arrow/remember}
	Esta flecha se guarda en la memoria del tablero, aunque el |goban| use |forget|.
\end{key}

\begin{key}{/pgfgo/arrow/forget}
	Esta flecha no se guarda en la memoria del tablero, aunque el |goban| use |remember|.
\end{key}

Las líneas tienen las mismas opciones, salvo |tip|:

\begin{key}{/pgfgo/line/color=\meta{color} (initially black)}
	Como |/pgfgo/arrow/color|.
\end{key}

\begin{key}{/pgfgo/line/width=\meta{dimension} (initially 1.2pt)}
	Como |/pgfgo/arrow/width|.
\end{key}

\begin{key}{/pgfgo/line/opacity=\meta{number} (initially 1)}
	Como |/pgfgo/arrow/opacity|.
\end{key}

\begin{key}{/pgfgo/line/halo=\meta{halo options} (default true)}
	Como |/pgfgo/arrow/halo|.
\end{key}

\begin{stylekey}{/pgfgo/line/halo/true}
	Como |/pgfgo/arrow/halo/true|.
\end{stylekey}

\begin{stylekey}{/pgfgo/line/halo/false}
	Como |/pgfgo/arrow/halo/false|.
\end{stylekey}

\begin{key}{/pgfgo/line/halo/color=\meta{color}}
	Como |/pgfgo/arrow/halo/color|.
\end{key}

\begin{key}{/pgfgo/line/halo/width=\meta{dimension} (initially 0.8pt)}
	Como |/pgfgo/arrow/halo/width|.
\end{key}

\begin{key}{/pgfgo/line/remember}
	Como |/pgfgo/arrow/remember|.
\end{key}

\begin{key}{/pgfgo/line/forget}
	Como |/pgfgo/arrow/forget|.
\end{key}
